% Options for packages loaded elsewhere
\PassOptionsToPackage{unicode}{hyperref}
\PassOptionsToPackage{hyphens}{url}
%
\documentclass[
  11pt,
]{article}
\usepackage{amsmath,amssymb}
\usepackage{iftex}
\ifPDFTeX
  \usepackage[T1]{fontenc}
  \usepackage[utf8]{inputenc}
  \usepackage{textcomp} % provide euro and other symbols
\else % if luatex or xetex
  \usepackage{unicode-math} % this also loads fontspec
  \defaultfontfeatures{Scale=MatchLowercase}
  \defaultfontfeatures[\rmfamily]{Ligatures=TeX,Scale=1}
\fi
\usepackage{lmodern}
\ifPDFTeX\else
  % xetex/luatex font selection
\fi
% Use upquote if available, for straight quotes in verbatim environments
\IfFileExists{upquote.sty}{\usepackage{upquote}}{}
\IfFileExists{microtype.sty}{% use microtype if available
  \usepackage[]{microtype}
  \UseMicrotypeSet[protrusion]{basicmath} % disable protrusion for tt fonts
}{}
\makeatletter
\@ifundefined{KOMAClassName}{% if non-KOMA class
  \IfFileExists{parskip.sty}{%
    \usepackage{parskip}
  }{% else
    \setlength{\parindent}{0pt}
    \setlength{\parskip}{6pt plus 2pt minus 1pt}}
}{% if KOMA class
  \KOMAoptions{parskip=half}}
\makeatother
\usepackage{xcolor}
\usepackage[margin=25mm]{geometry}
\usepackage{color}
\usepackage{fancyvrb}
\newcommand{\VerbBar}{|}
\newcommand{\VERB}{\Verb[commandchars=\\\{\}]}
\DefineVerbatimEnvironment{Highlighting}{Verbatim}{commandchars=\\\{\}}
% Add ',fontsize=\small' for more characters per line
\newenvironment{Shaded}{}{}
\newcommand{\AlertTok}[1]{\textcolor[rgb]{1.00,0.00,0.00}{\textbf{#1}}}
\newcommand{\AnnotationTok}[1]{\textcolor[rgb]{0.38,0.63,0.69}{\textbf{\textit{#1}}}}
\newcommand{\AttributeTok}[1]{\textcolor[rgb]{0.49,0.56,0.16}{#1}}
\newcommand{\BaseNTok}[1]{\textcolor[rgb]{0.25,0.63,0.44}{#1}}
\newcommand{\BuiltInTok}[1]{\textcolor[rgb]{0.00,0.50,0.00}{#1}}
\newcommand{\CharTok}[1]{\textcolor[rgb]{0.25,0.44,0.63}{#1}}
\newcommand{\CommentTok}[1]{\textcolor[rgb]{0.38,0.63,0.69}{\textit{#1}}}
\newcommand{\CommentVarTok}[1]{\textcolor[rgb]{0.38,0.63,0.69}{\textbf{\textit{#1}}}}
\newcommand{\ConstantTok}[1]{\textcolor[rgb]{0.53,0.00,0.00}{#1}}
\newcommand{\ControlFlowTok}[1]{\textcolor[rgb]{0.00,0.44,0.13}{\textbf{#1}}}
\newcommand{\DataTypeTok}[1]{\textcolor[rgb]{0.56,0.13,0.00}{#1}}
\newcommand{\DecValTok}[1]{\textcolor[rgb]{0.25,0.63,0.44}{#1}}
\newcommand{\DocumentationTok}[1]{\textcolor[rgb]{0.73,0.13,0.13}{\textit{#1}}}
\newcommand{\ErrorTok}[1]{\textcolor[rgb]{1.00,0.00,0.00}{\textbf{#1}}}
\newcommand{\ExtensionTok}[1]{#1}
\newcommand{\FloatTok}[1]{\textcolor[rgb]{0.25,0.63,0.44}{#1}}
\newcommand{\FunctionTok}[1]{\textcolor[rgb]{0.02,0.16,0.49}{#1}}
\newcommand{\ImportTok}[1]{\textcolor[rgb]{0.00,0.50,0.00}{\textbf{#1}}}
\newcommand{\InformationTok}[1]{\textcolor[rgb]{0.38,0.63,0.69}{\textbf{\textit{#1}}}}
\newcommand{\KeywordTok}[1]{\textcolor[rgb]{0.00,0.44,0.13}{\textbf{#1}}}
\newcommand{\NormalTok}[1]{#1}
\newcommand{\OperatorTok}[1]{\textcolor[rgb]{0.40,0.40,0.40}{#1}}
\newcommand{\OtherTok}[1]{\textcolor[rgb]{0.00,0.44,0.13}{#1}}
\newcommand{\PreprocessorTok}[1]{\textcolor[rgb]{0.74,0.48,0.00}{#1}}
\newcommand{\RegionMarkerTok}[1]{#1}
\newcommand{\SpecialCharTok}[1]{\textcolor[rgb]{0.25,0.44,0.63}{#1}}
\newcommand{\SpecialStringTok}[1]{\textcolor[rgb]{0.73,0.40,0.53}{#1}}
\newcommand{\StringTok}[1]{\textcolor[rgb]{0.25,0.44,0.63}{#1}}
\newcommand{\VariableTok}[1]{\textcolor[rgb]{0.10,0.09,0.49}{#1}}
\newcommand{\VerbatimStringTok}[1]{\textcolor[rgb]{0.25,0.44,0.63}{#1}}
\newcommand{\WarningTok}[1]{\textcolor[rgb]{0.38,0.63,0.69}{\textbf{\textit{#1}}}}
\setlength{\emergencystretch}{3em} % prevent overfull lines
\providecommand{\tightlist}{%
  \setlength{\itemsep}{0pt}\setlength{\parskip}{0pt}}
\setcounter{secnumdepth}{-\maxdimen} % remove section numbering
\ifLuaTeX
  \usepackage{selnolig}  % disable illegal ligatures
\fi
\usepackage{bookmark}
\IfFileExists{xurl.sty}{\usepackage{xurl}}{} % add URL line breaks if available
\urlstyle{same}
\hypersetup{
  pdftitle={P16.4: actual target estimates, generic incidence points, and a local formal bridge},
  pdfauthor={Unverified research continuation relative to six accepted caller grants},
  hidelinks,
  pdfcreator={LaTeX via pandoc}}

\title{P16.4: actual target estimates, generic incidence points, and a
local formal bridge}
\author{Unverified research continuation relative to six accepted caller
grants}
\date{7 October 2026}

\begin{document}
\maketitle

\section{Status and retained original
problem}\label{status-and-retained-original-problem}

Original P16.4 remains open in this submission. The new results are the
target lemmas and estimates numbered 1--6, the generic-incidence
construction with explicit initial point hypotheses numbered 7--10, and
the local Cartier and section-lifting propositions 12--13. The finite
graph theorem 11 is independently reconstructed. These results do not
constitute an unconditional all-genus curve equation or an
arbitrary-dimensional existence-and-uniqueness theorem. In particular,
none of the unreviewed ROOT observations or native191 conclusions is a
mathematical prerequisite. The submission has \texttt{candidate=false}.

The unchanged original domain is a complete characteristic-zero,
residue-characteristic-zero, real-valued non-Archimedean field \(K\), a
smooth geometrically integral projective \(K\)-variety \(X\) of
arbitrary dimension \(n\), and its actual ample rational line bundle
\(L\). Trivial, nontrivial, discrete, nondiscrete, and
non-algebraically-closed fields remain included. Write
\(T=X^{\mathrm{an}}\) and \(V=(L^n)>0\). The target is an arbitrary
positive Radon measure of raw mass \(V\), with the intrinsic
finite-energy condition that it integrate every actual finite-energy
potential. All psh metrics mean actual all-point decreasing
Fubini--Study net endpoints on this bundle; literal values \(-\infty\)
and supported analytic points are retained. A scalar potential \(u\)
represents the absolute metric \(\psi+u\) for a continuous psh reference
\(\psi\).

For continuous psh potentials the raw relative energy is \[
 E(u)=\frac1{n+1}\sum_{j=0}^n\int_T u\,\operatorname{MA}(u^{[j]},0^{[n-j]}).
\] On an actual singular endpoint its extension is \[
 E(u)=\inf\{E(v):v\geq u,\ v\text{ continuous actual psh}\},\qquad
 \mathcal E^1=\{u:E(u)>-\infty\}.
\] The operator used throughout is the genuine operator granted in O,
with mass \(V\). No inverse target construction is substituted for that
operator. For a raw-mass target set \[
 J(\mu)=\sup_{v\text{ continuous actual psh}}
              \left(E(v)-\int_T v\,d\mu\right).
\] Constants cancel because \(E(v+c)=E(v)+cV\). P16.3, P16.5, and P16.7
are not altered. Dimension zero is the one-point case with \(V=1\); all
new substantive arguments below have \(n\geq1\).

\section{Exact mathematical
prerequisites}\label{exact-mathematical-prerequisites}

The six accepted caller grants, including every supplied clarification,
are used only in their stated scopes. Their full bodies are extracted
verbatim from the supplied assignment into \texttt{sources/grant-O.tex},
\texttt{grant-D.tex}, \texttt{grant-N.tex}, \texttt{grant-T.tex},
\texttt{grant-Q.tex}, and \texttt{grant-C.tex}; the clarification files
accompany them. \texttt{input.json} preserves the complete input. These
mathematical labels are not registered note IDs. Local note n1 records
their allowed use; supporting note lists contain only emitted local IDs.

O supplies the actual positive mixed \(\mathcal E^1\) operator, all-test
integrability against outputs, weighted decreasing continuity on
arbitrary directed nets, extended energy continuity including
\(-\infty\), concavity, translation and monotonicity, admissible convex
curvature calculus, actual segment differentiation, reference
covariance, and complete-field pushforward. Its normalization
clarification is retained; we never assume a nonpositive endpoint has
nonpositive defining approximants. D supplies invariant actual descent,
extended energy equality, invariant Radon lift, and continuous-dual
identity for the completed algebraic closure. It supplies neither
arbitrary noninvariant singular-test descent nor general equation
existence. N requires an already actual continuous single-atom norming
reference; coercivity is not attainment. T covers the full ample
trivial-field equation. Q covers the split projective-line equation and
its local disk geometry. C supplies strict locality and comparison. No
additional equation, compactness, envelope, or uniqueness result is
inferred from these labels.

The proofs in Sections 3--6 use O and elementary compact Radon theory.
The generic-incidence proof uses explicitly stated actual point data and
ordinary algebraic foundations, with a separate cohomological hypothesis
for its independent complete-intersection proof. The node proof uses an
explicitly specified actual split semistable cover. The section-lifting
proof states the exact cohomology, completeness, and analytic/algebraic
section-identification premises it needs. Their applicability to every
original model is not asserted.

\section{Directed continuous Radon
integration}\label{directed-continuous-radon-integration}

\textbf{Lemma 1.} Let \(\lambda\) be a finite Radon measure on a compact
Hausdorff space. If \(g_\alpha\) is a directed increasing net of
nonnegative continuous functions and \(g=\sup_\alpha g_\alpha\), then \[
 \int g\,d\lambda=\sup_\alpha\int g_\alpha\,d\lambda.
\]

\textbf{Proof.} For each \(t\geq0\), \(U_t=\{g>t\}\) is the union of the
directed open sets \(U_{\alpha,t}=\{g_\alpha>t\}\). Any compact subset
of \(U_t\) has a finite subcover by these open sets; directedness puts
that finite subcover inside one \(U_{\alpha,t}\). Inner regularity gives
\(\lambda(U_t)=\sup_\alpha\lambda(U_{\alpha,t})\). For a finite height
cutoff and a finite lower step approximation to the layer integral,
choose a single index that approximates all its finitely many open-level
masses. Its integral is at least the resulting step sum. Refining the
steps and increasing the cutoff proves the reverse inequality to the
immediate monotone bound. This proves the formula including an infinite
integral. It is not a monotone-convergence assertion for arbitrary
measurable nets. \(\square\)

\section{Finite continuous dual implies the intrinsic
condition}\label{finite-continuous-dual-implies-the-intrinsic-condition}

\textbf{Theorem 2.} In the full original domain, \(J(\mu)<\infty\)
implies that every actual \(u\in\mathcal E^1\) is absolutely
\(\mu\)-integrable and that its pole set has zero \(\mu\)-mass. The
supremum of the objective over actual \(\mathcal E^1\) equals
\(J(\mu)\).

\textbf{Proof.} Choose a defining continuous FS net
\(f_\alpha\downarrow u\). Restrict to a tail dominated by one continuous
stage and hence by a finite constant \(M\). The nonnegative continuous
functions \(M-f_\alpha\) increase to the lower semicontinuous function
\(M-u\). By the dual inequality and O's energy monotonicity, \[
 \int_T(M-f_\alpha)\,d\mu
 \leq MV+J(\mu)-E(f_\alpha)
 \leq MV+J(\mu)-E(u)<\infty.
\] Lemma 1 gives the same finite bound for \(\int(M-u)d\mu\). Since
\(u\) is bounded above, this proves absolute integrability and the pole
assertion. The same lemma now gives convergence of the integrals. O's
extended decreasing-energy continuity gives
\(E(f_\alpha)\downarrow E(u)\). Consequently \[
 E(u)-\int_Tu\,d\mu
   =\lim_\alpha\left(E(f_\alpha)-\int_Tf_\alpha\,d\mu\right)\leq J(\mu).
\] Continuous tests belong to \(\mathcal E^1\), giving the reverse
supremum inequality. No norming point, normalized-family compactness, or
envelope theorem is used. \(\square\)

\section{A unit-energy bound without an
anchor}\label{a-unit-energy-bound-without-an-anchor}

Normalize an actual finite-energy potential by its finite supremum, so
\(u\leq0\) and \(\sup_Tu=0\). Put \[
 B(u)=-E(u)\geq0,\qquad I_\mu(u)=\int_T(-u)\,d\mu,
\] initially allowing \(I_\mu(u)=\infty\), and set \[
 U_\mu=\sup\{I_\mu(u):u\in\mathcal E^1,\ \sup_Tu=0,\ B(u)\leq1\}.
\]

\textbf{Theorem 3.} In the full original domain, integrability of every
actual \(\mathcal E^1\) potential is equivalent to \(U_\mu<\infty\). If
\(U_\mu\) is finite, then every normalized \(u\) satisfies \[
 I_\mu(u)\leq U_\mu\max\{1,B(u)\}.
\] This linear estimate alone does not imply finite continuous dual.

\textbf{Proof.} Suppose all actual tests are integrable but
\(U_\mu=\infty\). Choose normalized \(u_j\) with \(B(u_j)\leq1\) and
\(I_\mu(u_j)\geq4^j\). Let \[
 v_N=\sum_{j=1}^N2^{-j}u_j+2^{-N}0.
\] O's convexity and concavity give actual finite-energy \(v_N\),
\(v_{N+1}\leq v_N\leq0\), and \(E(v_N)\geq-\sum_{j=1}^N2^{-j}\geq-1\).
Let \(v=\inf_Nv_N\) at every point. This limit is not identically
\(-\infty\). Otherwise, for \(M>1/V\), the nested closed sets
\(\{v_N\geq-M\}\) would have empty intersection. Compactness gives an
empty member, hence \(v_N<-M\) everywhere and, by monotonicity,
\(E(v_N)\leq E(-M)=-MV<-1\), a contradiction. This is compactness of the
underlying analytic carrier, not compactness of a normalized psh family.

Here is an explicit actual FS recognition of this endpoint, avoiding a
diagonal sequence on a nonmetrizable space. Let \(\mathcal F\) contain
all continuous weighted FS scalars \(f\) for which
\(f\geq v_N+\epsilon\) for some \(N\) and \(\epsilon>0\). It is
nonempty: add \(\epsilon\) to a defining FS approximant of \(v_N\).
Every member majorizes \(v\). For any fixed \(N\), those shifted
approximants have infimum \(v_N\) as the shift tends to zero, so
\(\inf\mathcal F=v\) pointwise.

The family is downward directed. Given \(f\geq v_N+\epsilon\) and
\(g\geq v_M+\delta\), choose \(L\geq\max(N,M)\) and
\(0<\eta<\min(\epsilon,\delta)/4\). A defining continuous FS net
\(h_\alpha\downarrow v_L\) eventually satisfies
\(h_\alpha\leq\min(f,g)-2\eta\) on the whole compact carrier: its closed
upper-level sets against the continuous function \(\min(f,g)-2\eta\)
decrease to the empty set, because the endpoint has a strictly positive
gap. Then \(h_\alpha+\eta\) belongs to \(\mathcal F\) and lies below
both \(f,g\). Ordered by reverse pointwise inequality, \(\mathcal F\)
therefore defines an actual decreasing FS net to \(v\). Quotient equal
functions if a partial order is desired.

Thus \(v\) is a nontrivial actual psh endpoint. Every \(f\in\mathcal F\)
has \(E(f)\geq E(v_N+\epsilon)\geq-1\). O's \textbf{extended}
decreasing-energy continuity, applied before asserting \(\mathcal E^1\)
membership, gives \(E(v)=\inf_{\mathcal F}E(f)\geq-1\). Hence
\(v\in\mathcal E^1\). Ordinary sequential monotone convergence now gives
\[
 I_\mu(v)=\sum_{j=1}^\infty2^{-j}I_\mu(u_j)
          \geq\sum_{j=1}^\infty2^j=\infty,
\] contradicting the hypothesis.

Conversely, if \(U_\mu<\infty\), O's concavity gives \(B(tu)\leq tB(u)\)
for \(0<t\leq1\). Take \(t=1\) when \(B(u)\leq1\) and \(t=1/B(u)\)
otherwise. The scaled actual potential has supremum zero and cost at
most one. Nonnegative integration gives
\(tI_\mu(u)=I_\mu(tu)\leq U_\mu\), proving the estimate and finiteness.
Other potentials differ by a finite constant. \(\square\)

\section{The exact scaling bound and the missing geometric
estimate}\label{the-exact-scaling-bound-and-the-missing-geometric-estimate}

\textbf{Proposition 4.} For every normalized actual
\(u\in\mathcal E^1\), set \[
 A_k(u)=\int_T(-u)\,\operatorname{MA}(u^{[k]},0^{[n-k]}),\qquad
 C_n=\frac{n(n+1)}2.
\] Then all \(A_k\) are finite nonnegative,
\(B=(n+1)^{-1}\sum_{k=0}^nA_k\), and \[
 B(tu)\leq tA_0(u)+C_nt^2B(u)\qquad(0\leq t\leq1).
\]

\textbf{Proof.} O supplies the finite integrals, genuine energy
identity, and affine curvature calculus. For \(0<t\leq1\) their finite
positive expansion is \[
 B(tu)=\frac t{n+1}\sum_{j=0}^n\sum_{k=0}^j
       \binom jk t^k(1-t)^{j-k}A_k(u).
\] The \(k=0\) terms sum to at most \(tA_0\). For \(k\geq1\),
\(A_k\leq(n+1)B\) and the binomial weight in the \(j\)th sum is
\(1-(1-t)^j\leq jt\). Their total is at most
\(tB\sum_{j=0}^n jt=C_nt^2B\). At \(t=0\) use the reference potential
itself; no zero multiple of a pole is evaluated. \(\square\)

\textbf{Corollary 5, conditional on a uniform reference bound.} Suppose,
additionally, that there is a proved finite \(M_0\) with
\(A_0(u)\leq M_0\) for every normalized actual finite-energy potential
on the reference and bundle under consideration. Then the intrinsic
condition, \(U_\mu<\infty\), and \(J(\mu)<\infty\) are equivalent. With
\[
 c=2(M_0+1+\sqrt{C_n}),\qquad A=cU_\mu,
\] one has \[
 I_\mu(u)\leq A(1+\sqrt{B(u)}),\qquad
 E(u)-\int u\,d\mu\leq-B(u)+A(1+\sqrt{B(u)}),\qquad
 J(\mu)\leq A+\frac{A^2}4.
\]

\textbf{Proof.} Take \(t=[c(1+\sqrt{B(u)})]^{-1}\). Then \(0<t\leq1\),
\(tM_0\leq1/2\), and \(C_nt^2B(u)\leq1/4\). Proposition 4 gives
\(B(tu)\leq3/4\). Theorem 3's unit bound gives \(tI_\mu(u)\leq U_\mu\)
and the displayed estimates. Maximizing \(-b+A\sqrt b\) over \(b\geq0\)
gives \(A^2/4\). Theorem 2 and Theorem 3 supply the other implications.
The objective tends to \(-\infty\) when \(B\to\infty\). This proves
coercivity, not attainment. The required \(M_0\) has \textbf{not} been
proved for all original varieties here. \(\square\)

There is an unconditional necessary-target conclusion which does not
require \(M_0\).

\textbf{Corollary 6.} Every genuine output \(\mu=\operatorname{MA}(w)\),
\(w\in\mathcal E^1\), has finite continuous dual, and \(w\) attains the
objective supremum over all actual \(\mathcal E^1\): \[
 J(\mu)=E(w)-\int_Tw\,d\mu<\infty.
\]

\textbf{Proof.} O integrates every actual test against this output.
Concavity of the actual segment energy and its derivative at \(w\) give
\[
 E(v)-E(w)\leq\int_T(v-w)\,d\mu
\] for every actual \(v\in\mathcal E^1\). Each integral is finite on a
common finite-value full-measure set; common pole values are never
subtracted. This bounds every objective by the displayed finite value.
Continuous majorants \(w_\alpha\downarrow w\) converge in energy and in
their integrals against the fixed output by O's weighted law, so the
continuous supremum has this value as well. This proves a necessary
condition and an attainment statement for already known outputs. It does
not construct an output for an arbitrary prescribed intrinsic target.
\(\square\)

\section{Actual generic incidence: hypotheses and common
field}\label{actual-generic-incidence-hypotheses-and-common-field}

The following construction is separate from any unproved Shilov-model
assertion. Let \(C\) be complete, algebraically closed, nontrivially
valued, of characteristic and residue characteristic zero. Let \(X/C\)
be smooth geometrically integral projective of dimension \(n\geq2\).
Clear rational denominators and choose an actual very ample \(M=mL\),
with an embedding given by \(S_0,\ldots,S_N\). Put \(g_l=S_l/S_0\) on
the affine chart \(S_0\ne0\).

Assume finitely many actual points \(y_i\) have generic scheme support
and satisfy the following explicit point hypotheses:

\begin{itemize}
\tightlist
\item
  \(|g_l(y_i)|\leq1\) for every \(l\);
\item
  the reductions of the \(g_l(y_i)\) generate a subfield
  \(k(Y_i)\subset\widetilde{\mathcal H(y_i)}\), of transcendence degree
  \(n\) over \(k=\widetilde C\), and are coordinates of an integral
  affine \(k\)-variety \(Y_i\);
\item
  \(\operatorname{trdeg}_k\widetilde{\mathcal H(y_i)}=n\).
\end{itemize}

Equality of \(k(Y_i)\) with the whole completed residue field's residue
is not required. These are initial hypotheses. This paper does not prove
they hold for every Shilov point of a suitable actual model reference.

Set \(s=n-1\), \(r'=sN\), and \(r=s(N+1)\). Let \(H_i=\mathcal H(y_i)\),
the completion of the honest valued field \(C(X)\). Introduce \(r'\)
free variables \(A_{jl}\) (\(1\leq j\leq s\), \(1\leq l\leq N\)) with
unit-radius Gauss norm over \(H_i\), and complete to \(H'_i\). Define in
this field \[
 A_{j0}=-\sum_{l=1}^N A_{jl}g_l(y_i).
\]

\textbf{Lemma 7, residual independence.} All \(r\) parameters, including
the solved \(A_{j0}\), have algebraically independent residues over
\(k\).

\textbf{Proof.} Consider the algebraic map \[
 Y_i\times\mathbb A^{r'}_k\longrightarrow\mathbb A^r_k,
 \quad (y,a)\longmapsto\bigl(a_{jl},-\sum_l a_{jl}\bar g_l(y)\bigr).
\] The integral variety \(Y_i\) has a nonempty smooth locus in
characteristic zero. The coordinate functions \(\bar g_l\) give its
closed affine embedding, so at a smooth point their differentials span
its \(n\)-dimensional cotangent space. Choose \(s\) independent linear
combinations of these differentials. At the corresponding coefficient
rows, the differential of the displayed map has the identity block on
all \(r'\) free coordinates and an additional rank-\(s\) block in the
\(Y_i\) directions. Hence an \(r\)-minor is nonzero at that point, and
at the generic point. In characteristic zero, \(r\) independent
differentials of functions imply their algebraic independence: the rank
of the differentials of a subfield equals its transcendence degree, by
choosing a transcendence basis and using separability of algebraic
extensions. Thus the map is dominant and the \(r\) output functions are
algebraically independent in \(k(Y_i)(A_{\mathrm{free}})\). Its
injection into \(\widetilde H_i(A_{\mathrm{free}})\) preserves that
assertion. \(\square\)

\textbf{Lemma 8, the same complete coefficient-Gauss field.} Let \[
 F=\widehat{C(A_{jl}:1\leq j\leq s,\ 0\leq l\leq N)}
\] with the unit coefficient-Gauss norm on all \(r\) independent
symbols. The preceding evaluation embeds this \textbf{same} complete
field \(F\) isometrically in every \(H'_i\).

\textbf{Proof.} Each full parameter is integral and its residue is a
nonzero member of an independent family, hence its norm is one. For any
nonzero polynomial over \(C\), divide by a coefficient of maximal norm.
Its reduction is a nonzero polynomial over \(k\) and cannot vanish on
the independent residues from Lemma 7. The evaluation therefore has norm
exactly one after normalization, proving the coefficient-Gauss identity.
Fractions retain their quotient norms, so the rational function field
map is injective and isometric. Completeness extends it uniquely and
isometrically to \(F\). The formulas representing the solved
coefficients can differ with \(i\); the domain field and its norm are
identical. Completion of a valued field changes neither its residue
field nor its value group: approximating a nonzero completed element to
strictly smaller error preserves its norm and, after unit normalization,
its residue. In particular \(\widetilde F=k(A_{\mathrm{all}})\) and
\(\widetilde H'_i=\widetilde H_i(A_{\mathrm{free}})\). \(\square\)

\section{A complete-intersection proof with an explicit cohomological
premise}\label{a-complete-intersection-proof-with-an-explicit-cohomological-premise}

For an independent proof of geometric integrality, impose the following
finite cohomological hypothesis on the chosen actual embedding power: \[
 H^q(X,\mathcal O_X(-kM))=0
 \quad(1\leq k\leq n-1,\ 1\leq q\leq n-1). \tag{H}
\] This is an explicit premise of the construction as formulated here.
Smooth-projective Serre duality and Serre vanishing give such
sufficiently large powers; their complete recursive proofs are not newly
supplied in this paper. No unverified model cohomology vanishing over
\(C^\circ/a\) is substituted for (H). Alternatively, the complete
generic Bertini primary units in the ledger give a source route to
geometric irreducibility. We record their foundations separately instead
of declaring them new caller grants.

\textbf{Proposition 9.} Under (H), the \(s=n-1\) full generic
hyperplanes \[
 A_{j0}S_0+\sum_{l=1}^NA_{jl}S_l=0\quad(1\leq j\leq s)
\] cut out a smooth geometrically integral projective curve \(D\) over
\(F\). Its actual restricted bundle has raw degree \[
 a=\deg(L|_D)=m^{n-1}(L^n).
\]

\textbf{Proof.} First work over the rational coefficient field
\(E=C(A_{\mathrm{all}})\). Here the coefficients are the full
independent generic coefficients, not a fixed projection subsystem. For
a smooth embedded projective variety \(Z\) of dimension \(d\), the
conditions that a hyperplane section vanish and have zero differential
at a specified point impose \(d+1\) independent linear conditions on its
coefficients. Indeed choose a nonvanishing embedding section as local
frame; the resulting affine coordinate functions span the cotangent
space. The incidence of such bad pairs therefore has dimension \(N-1\)
over the parameter space of projective hyperplanes of dimension \(N\).
Its projection is closed because \(Z\) is projective. It misses the
generic hyperplane. The calculation also holds after algebraic closure.
Thus the generic section is geometrically smooth. On each
positive-dimensional component the generic hyperplane neither contains
the component nor misses it, by the projective intersection dimension
theorem; its intersection has dimension \(d-1\).

Apply this argument successively with the independent rows of full
coefficients. Each intermediate intersection is smooth, has the expected
dimension, and the next equation is nonzero on every component. The
equations form a regular sequence; its Koszul complex gives a resolution
of the final structure sheaf by finite direct sums of
\(\mathcal O_X(-kM)\), \(1\leq k\leq s\). For clarity, the degree-zero
hypercohomology of this resolution has only \(H^0(\mathcal O_X)\)
contributing: a term \(\mathcal O_X(-kM)\) lies in degree \(-k\) and
would require \(H^k\) to contribute to degree zero, or \(H^{k-1}\) to
map into it. These vanish by (H) and \(H^0\) of a negative ample bundle
is zero. Therefore \(H^0(D_E,\mathcal O_{D_E})=E\). Cohomology commutes
with extension of the base field: a finite affine Čech cover of a
separated scheme has affine intersections, and tensoring its complex
with a field is exact. The same vanishings and equality hold over an
algebraic closure and over the completion \(F\).

The geometrically smooth projective final curve has no nilpotents, and
different irreducible components of a smooth curve are disjoint. If
there were two geometric components, they would give a nontrivial
idempotent in its global functions. The equality just proved excludes
this. It is therefore geometrically integral. The same Čech and
geometric base-extension argument preserves this conclusion over \(F\).
Each equation is an actual section of the actual \(M\) on \(X_F\), so
the regular-sequence intersection class is \(M^{n-1}\), and intersecting
it with \(L\) gives the displayed raw degree. Rational powers were
cleared before defining this embedding. \(\square\)

The dimension and cotangent foundations in this proof are ordinary
algebraic foundations, not a proof of the formal model assertions
R1--R5. The cohomological premise makes explicit where a sufficiently
high embedding power is used. In particular the generic curve is not
claimed to be genus zero.

\section{Actual incidence points and exact preservation of original
values}\label{actual-incidence-points-and-exact-preservation-of-original-values}

\textbf{Theorem 10.} Under the preceding point hypotheses and (H), there
is an actual analytic point \(y'_i\in D^{\mathrm{an}}\) over each
\(y_i\), with \[
 \mathcal H(y'_i)=H'_i,\qquad
 \operatorname{trdeg}_{\widetilde F}\widetilde{\mathcal H(y'_i)}=1,
 \qquad |\mathcal H(y'_i)^\times|=|F^\times|=|C^\times|.
\] Every original rational-function value, and every original section
norm in a pulled-back continuous metric, is preserved. No point of a
curve over the unchanged ground field is asserted.

\textbf{Proof.} In the affine chart the map
\(F[g_1,\ldots,g_N]\to H'_i\) given by the original evaluations kills
all the hyperplane equations, by their defining solved coefficients.
Hence it defines a multiplicative seminorm on the actual coordinate
algebra of \(D\); this is an analytic point. Restricted to the original
\(C\)-coordinate algebra it is exactly \(y_i\).

The kernel on the integral curve algebra is zero. To see this without
assuming a function-field injection in advance, its seminorm residue
field contains the reductions of the original \(g_l\) and the free
parameters. These generate \(k(Y_i)(A_{\mathrm{free}})\), of
transcendence degree \(n+r'\) over \(k\). The base residue field is
\(\widetilde F=k(A_{\mathrm{all}})\), of transcendence degree
\(r=r'+n-1\); it lies in that same field by the incidence formulas. Thus
the point's residue has relative transcendence degree at least one. A
nongeneric point of an integral affine \(F\)-curve is closed and has
finite residue extension of \(F\). The valuation residue of a finite
valued extension is algebraic: lifts of a residue-field linearly
independent finite set are linearly independent over the base field,
after normalizing a relation by a maximal coefficient. Its residual
transcendence degree is zero. This contradiction proves generic support
and injection of \(F(D)\).

The image of \(F(D)\) contains \(C(X)\) and all free coefficient
variables. Indeed the \(g_l\) generate the original affine coordinate
ring's function field. Since \(y_i\) has generic support, \(C(X)\) is
dense in \(H_i\) by the definition of its completed residue field.
Polynomials in the free variables with coefficients in \(H_i\) are
Gauss-approximated by polynomials with coefficients in \(C(X)\), by
approximating their finitely many coefficients. For a nonzero
denominator, a sufficiently small approximation error preserves its norm
and controls its inverse. Thus \(C(X)(A_{\mathrm{free}})\) is dense in
\(H_i(A_{\mathrm{free}})\) and in its completion \(H'_i\). The completed
residue field of the point, being the closure of \(F(D)\) in \(H'_i\),
is exactly \(H'_i\).

The residue formula in Lemma 8 and the initial maximal residual
dimension give \[
 \operatorname{trdeg}_{\widetilde F}\widetilde H'_i
    =(n+r')-(r'+n-1)=1.
\] There is no need for equality \(\widetilde H_i=k(Y_i)\). All free
Gauss variables have norm one. Here is the elementary valuation
independence bound used for the value group. Choose units with \(d\)
independent residues, and field elements with \(q\) rationally
independent values modulo the base value group. A polynomial in these
elements can be grouped by the exponents of the latter \(q\) elements.
Within each group, divide its coefficient polynomial in the units by a
base coefficient of maximal norm; independent residues make that
coefficient polynomial's norm exactly the maximal coefficient norm.
Different groups have different valuations, since their value
differences cannot lie in the base group. A unique least valuation
therefore survives in any nonzero polynomial. Thus the \(d+q\) chosen
elements are algebraically independent, and \(d+q\leq n\) in \(C(X)\).
Residue and group are unchanged by completion. Taking \(d=n\) shows the
value-group extension is torsion. Since \(C\) has divisible value group,
an ordered torsion extension inside \(\mathbb R\) of that subgroup has
no new element: if an integer multiple of a real value belongs to the
subgroup, division by that integer puts the value itself there.
Consequently \(|H_i^\times|=|C^\times|\), and the unit Gauss extension
and completion preserve this equality.

Finally, \(H_i\hookrightarrow H'_i\) is isometric. Evaluation of every
original rational function is unchanged. For an original section, use a
local rational frame nonvanishing at the generic support; its
coefficient norm is unchanged, and the pulled-back metric's frame norm
is its original value. This proves the section assertion. \(\square\)

There is also an actual completed-algebraic-closure lift preserving
these invariants. Put \(B=\widehat{\overline F}\). Extend
\(F\hookrightarrow H'_i\) to an embedding of an algebraic closure into
an algebraic closure of \(H'_i\); uniqueness of the extension of the
valuation from the complete field \(F\) makes it isometric. It extends
to \(B\hookrightarrow\Omega_i=\widehat{\overline{H'_i}}\). The original
coordinate evaluations now define a point \(z_i\in D_B^{\mathrm{an}}\).
The injection of the old \(F(D)\) survives evaluation, so its closure
\(H'_i\) embeds in \(\mathcal H(z_i)\). Residues of \(\Omega_i\) are
algebraic over \(\widetilde H'_i\), and \(\widetilde B\) is algebraic
over \(\widetilde F\): algebraic valuation extensions have algebraic
residue, and completion changes no residue. The residue of this point
thus has transcendence degree one over \(\widetilde B\). The same
closed-point argument as above proves its generic support before using
\(B(D)\). Its completed residue field is then exactly the closure of
\(B(D)\) in \(\Omega_i\). The value group stays the original divisible
group. These are the usual type-II invariants on an algebraically closed
curve. Recognition by a semistable model still uses the curve geometry
and formal/reduction foundations audited below. GM A.7 has \textbf{not}
been applied to the initial transcendental extension \(F/C\); its stated
domain is algebraic-completion extensions only.

Theorem 10 constructs the same coefficient-Gauss base at every selected
original point and preserves actual values. It does not prove the
initial Shilov hypotheses, semistable containment of the selected lifts,
the raw curve curvature formula for all already actual restricted
metrics, or the uniform \(M_0\) required in Corollary 5.

\section{Positive rational-slope interpolation with irrational
cycles}\label{positive-rational-slope-interpolation-with-irrational-cycles}

\textbf{Theorem 11, a finite graph theorem.} Let
\(\Lambda\subset\mathbb R\) be a nonzero divisible subgroup. Let \(G\)
be a finite connected metric graph with edge lengths in \(\Lambda\) and
fixed positive real weights \(w_e\). Let \(\theta\) be a signed vertex
measure and \(\mu\) a nonnegative vertex measure, both of mass \(V>0\).
For a vertex \(z\) there is a unique edge-affine \(F\) with \(F(z)=0\)
and \(\theta+\Delta F=\mu\), where \(\Delta\) is weighted outgoing
slope. There are subdivisions and rational-slope functions \(G_j\) with
all vertex values in \(\Lambda\), one common denominator at each stage,
\(G_j\to F\) uniformly, \(\theta+\Delta G_j\geq0\), and
\(\theta+\Delta G_j\to\mu\) weakly. Every vertex of the refined graph
has strictly positive curvature at each stage.

\textbf{Proof.} Subdivide every self-loop at its midpoint, retaining the
loop and giving the new vertex zero initial mass. For vertex values
\(h\), edgewise cancellation gives \[
 -\sum_vh(v)\Delta h(v)=\sum_{e=ab}\frac{w_e}{\ell_e}(h(b)-h(a))^2.
\] The matrix is symmetric with kernel exactly constants by
connectivity, hence image the zero-sum vectors. This gives the anchored
solution. A no-edge singleton has only the constant solution.

For the nontrivial graph, let \(k\) be its number of vertices and
replace \(\mu\) by
\(\mu_\epsilon=(1-\epsilon)\mu+\epsilon(V/k)\sum_v\delta_v\). Its
anchored solution \(F_\epsilon\) converges to \(F\), by
finite-dimensional linearity, and every old vertex has a positive
margin. Since \(\mathbb Q a\subset\Lambda\) is dense for any nonzero
\(a\in\Lambda\), approximate the finitely many shared vertex values by
\(g_v\in\Lambda\), with \(g_z=0\) and error at most \(h\). On an
oriented edge \(a\to b\), write \(d=g_b-g_a\), \(c=d/\ell\). Choose
rational \(q_1<c<q_2\) as close to \(c\) as desired. Set \[
 t=\frac{d-q_2\ell}{q_1-q_2}\in\Lambda,\qquad 0<t<\ell.
\] Use slope \(q_1\) on the first subedge and \(q_2\) on the second. The
terminal value is exactly \(g_b\), and the bend value \(g_a+q_1t\) lies
in \(\Lambda\). Shared endpoint values guarantee exact compatibility
around every cycle.

The bend contributes \(w_e(q_2-q_1)>0\). The two endpoint contributions
are \(w_eq_1\) and \(-w_eq_2\). Taking \(h\) small enough and both
rational slopes within \(\eta\) of the original \(F_\epsilon\) slope
makes the error at each old vertex less than \(\eta\sum_{e\ni v}w_e\).
Choose this below half its positive margin. Thus all old and new
curvature atoms are positive and their total mass remains \(V\). The
uniform function error is at most \(h+\eta\max_e\ell_e\); the inserted
mass tends to zero with \(\eta\). Letting the errors and \(\epsilon\)
tend to zero gives both convergences. A finite common denominator clears
every rational slope, with all scaled vertex values still in
\(\Lambda\). \(\square\)

The simpler fixed-edge rational-slope approach fails on two parallel
edges of lengths \(1\) and \(\sqrt2\): their shared endpoint difference
would have to lie in \(\mathbb Q\cap\sqrt2\mathbb Q=\{0\}\). The
positive movable bends are essential. For \(\theta=(1,1)\),
\(\mu=(2,0)\), \(\epsilon=1/10\), and \(d=(9/10)(2-\sqrt2)\), slopes
\(52/100,53/100\) and \(37/100,38/100\) give bend positions
\(-127+90\sqrt2\) and \(-180+128\sqrt2\). The curvature is
\(189/100,9/100,1/100,1/100\), with total exactly \(2\). The exact
program independently checks these formulas and 81 specified variants.
It supplies no actual formal model.

\section{A local actual vertical Cartier
calculation}\label{a-local-actual-vertical-cartier-calculation}

\textbf{Proposition 12, with explicit actual cover hypotheses.} Work
over \(R=C^\circ\) with \(C\) complete algebraically closed and
nontrivial. Suppose an actual proper split strict semistable formal
curve has a node-and-vertex cover of the following form. A node block
for an edge \(a\to b\) is étale over \[
 \operatorname{Spf}R\langle x,y\rangle/(xy-\pi_e),\qquad
 \ell_e=\operatorname{val}(\pi_e)>0.
\] The endpoint \(a\) has \(\operatorname{val}(x)=0\), and \(b\) has
\(\operatorname{val}(x)=\ell_e\). On intersections with the smooth
vertex part at \(a\), \(x\) is a formal unit; at \(b\), \(y\) is a
formal unit. Distinct node blocks meet through those smooth parts. The
resulting reductions identify the residual branches as the usual
parameters at their nodes, with no extra zeros or poles. Self-loops have
first been subdivided, so their two branches remain present. Let \(f\)
be integral \(\Lambda\)-affine on the skeleton and constant on
retraction fibres, with \(f(v)\in\Lambda=\operatorname{val}(C^\times)\).

Then \(f\) defines an actual vertical Cartier line bundle on this formal
model, generically trivialized by \(1\), of weight \(f\circ r\). Its
residual component degree at \(v\) is \(\Delta f(v)\) in the unweighted
outgoing-slope convention. Tensoring an already actual formal bundle
changes its actual residual degrees by these numbers and retains its
actual generic bundle. Rational slopes are handled by clearing one
common denominator.

\textbf{Proof.} Choose \(c_v\in C^\times\) with valuation \(f(v)\). On
the smooth vertex block use the meromorphic equation \(g_v=c_v\). On an
oriented node block, where \(f=f(a)+q\operatorname{val}(x)\) and
\(q\in\mathbb Z\), use \(g_e=c_a x^q\). The constants \(c_v\) and the
coordinates \(x,y\) are non-zero-divisors; \(xy=\pi_e\) and flatness
imply this for \(x,y\). Negative exponents lie in the total meromorphic
ring. They are all invertible on the generic fibre.

On the \(a\) overlap, \(g_e/g_a=x^q\) is a literal formal unit. On the
\(b\) overlap, \[
 g_e/g_b=(c_a\pi_e^q/c_b)y^{-q}
\] is a unit: the constant factor has valuation zero by
\(f(b)=f(a)+q\ell_e\), and \(y\) is a formal unit there. Overlaps of
node blocks reduce to these verified vertex ratios. They consequently
satisfy the cocycle identity as ratios in a common meromorphic ring.
This constructs the Cartier divisor \(D\), rather than merely a function
of norm one. In the local frame \(g^{-1}\) of \(\mathcal O(D)\), the
canonical generic section \(1\) has norm \(|g|\), so its additive weight
is \(\operatorname{val}(g)=f\circ r\).

For the degree on component \(C_a\), normalize its rational generic
section by \(1/c_a\). Near the node its coefficient in the local frame
is \(x^q\), so its reduction has order \(q\) at that branch. For
\(C_b\), normalization by \(1/c_b\) gives a unit times \(y^{-q}\), hence
order \(-q\). On the smooth vertex part the coefficient is one; there
are no other zeros or poles by the cover hypothesis. Summing the orders
of a rational section on the normalization of each proper residual
component gives its line-bundle degree. Thus the correction is the sum
of outgoing slopes. Each edge's two corrections sum to zero;
subdivisions and loops preserve the calculation. For \(f/m\), apply the
construction to \(f\) and take the rational tensor power \(1/m\).
\(\square\)

This calculation is confined to this explicitly given actual cover. It
does not prove global semistable domination, algebraization, the cover
theorem itself, or the genuine raw component formula for the analytic
Monge--Ampère operator. It avoids KRZB's analytic-meromorphic extension
only for this vertical-monomial degree calculation; that extension is
not declared unnecessary for other functions.

\section{Section lifting and actual pure FS
recognition}\label{section-lifting-and-actual-pure-fs-recognition}

\textbf{Proposition 13, with explicit cohomology and GAGA hypotheses.}
Let \(R=C^\circ\) and \(0<|a|<1\). Suppose \(\mathfrak X\) is a proper
flat, \(a\)-adically complete, separated formal \(R\)-curve, locally
with complete separated rings, and \(\mathfrak A\) an actual formal
invertible sheaf. Its generic analytification is identified with the
original \(X^{\mathrm{an}}\), and \(\mathfrak A\) there is identified
with the actual \(dL\) after clearing denominators. Assume for some
\(k\geq1\):

\begin{enumerate}
\def\labelenumi{\arabic{enumi}.}
\tightlist
\item
  \(\mathfrak A_0^k\) on \(X_0=\mathfrak X\bmod a\) is generated by
  finitely many global sections and \(H^1(X_0,\mathfrak A_0^k)=0\);
\item
  global analytic sections of its generic fibre coincide, under the
  given bundle identification, with algebraic sections of the actual
  \(kdL\) on \(X\).
\end{enumerate}

Then the formal metric on \(dL\) is an actual pure Fubini--Study metric,
with the same bundle and no reference to a numerical replacement.

\textbf{Proof.} Put \(A_r=\mathfrak A^k/a^{r+1}\mathfrak A^k\). Flatness
and local freeness give identical associated-graded layers
\(a^r\mathfrak A^k/a^{r+1}\mathfrak A^k\simeq A_0\). Thus for every
\(r\geq1\) there is an exact sequence of sheaves \[
 0\longrightarrow A_0\longrightarrow A_r\longrightarrow A_{r-1}\longrightarrow0.
\] The single stated \(H^1\) vanishing makes
\(H^0(A_r)\to H^0(A_{r-1})\) surjective for every \(r\). Lift each of
the finitely many chosen generating sections through all these maps,
compatibly. Sheaf completeness identifies global sections of
\(\mathfrak A^k\) with the inverse limit of these sections; this follows
from the defining complete local sheaves, since global sections commute
with limits. The lifted sections generate the formal bundle. Locally,
mod-\(a\) generation allows a chosen frame coefficient \(b\) to be
written \(\sum b_if_i+a b^{(1)}\). Repeating this on the successive
error, with all coefficients integral, produces convergent \(a\)-adic
sums of the \(b_i\). Completeness gives \(b=\sum b_if_i\) exactly. Hence
the local generation assertion does not assume the ring is Noetherian.

On the generic fibre each lifted section is algebraic by assumption 2.
They have no common zero by the formal generation. In a formal local
frame the coefficient functions satisfy \(|f_i|\leq1\). Integral
coefficients in an identity \(1=\sum b_if_i\) also have norm at most
one, so at every analytic point \(\max_i|f_i|=1\). The formal metric on
\(\mathfrak A^k\) therefore equals the metric defined by precisely these
global sections with zero FS weights. Dividing by \(k\) gives the metric
on the actual \(dL\), and by \(d\) the rational metric on \(L\). All
identities hold at every point of the original carrier. \(\square\)

This lemma needs one vanishing on the \textbf{whole mod-\(a\) fibre},
not only on the reduced residual curve. It also needs the actual
generic-bundle identification and the stated analytic/algebraic section
equality. Those inputs are not proved for every semistable model in this
paper. If the mod-\(a\) sections instead define a closed embedding and
their standard-chart preimages are affine formal schemes, the same
complete lifting argument makes the maps from completed projective chart
rings surjective: lift a polynomial representative modulo \(a\),
subtract it, divide by \(a\), and iterate. This gives a closed formal
embedding. Algebraizing its coherent ideal still requires projective
GFGA's coherence, full-faithfulness, and generation premises; this
observation is not an algebraization theorem.

The potential difficulty is concrete. For a dense divisible value group
choose \(\operatorname{val}(a)=1\). For every \(N\) there is \(b_N\)
with \(\operatorname{val}(b_N)=1/(N+1)\). Its image in \(R/aR\) is
nilpotent, but \(b_N^N\notin aR\). Thus the nilradical of \(R/aR\) is
not nilpotent. One may not justify whole mod-\(a\) ampleness or uniform
Serre vanishing by treating this fibre as a finite nilpotent thickening.
A separate valid argument, for example a fully verified finitely
presented Artinian descent, is needed.

\section{What the accessible primaries actually
resolve}\label{what-the-accessible-primaries-actually-resolve}

All 16 exact frozen primary files match their supplied bytes and SHA-256
hashes. All six extracted grant files also match the review-source
hashes supplied with their clarifications. These checks verify
provenance, not theorems. The quote-free ledger records whole printed
proof units actually read, scan checks, application scopes, and
recursive gaps. No prior retrieval failure is used to avoid an
accessible primary.

The most relevant source facts and boundaries are as follows.

\begin{itemize}
\tightlist
\item
  BL85's complete Section 7 was read, including the two
  genus-at-least-one lemmas, the induction and the final algebraic
  reduction. The scanned pages 374--378 were inspected. Its theorem has
  a stable genus-at-least-two branch and a semistable genus-at-most-one
  branch. The affinoid \(\mathbb P^1\) subdivision in its covering
  argument is explicitly omitted and refers to BGR. Earlier whole Lemma
  3.5 and relevant preceding annular arguments were read and checked
  against scans. Their earlier reduction, genus, finite-map, and
  formal-fibre dependencies are not all reconstructed. An incomplete
  recursive audit is not an unconditional new semistable theorem.
\item
  BPR's entire Theorem 4.11 proof 4.15 was read, including the one-edge
  and singleton cases, its supporting Lemmas 4.12 and 4.14, and its
  omitted analogous domination cases. Proposition 3.13 and the
  inverse-limit proof 5.2 were read. The latter invokes a separation
  assertion and compact inverse systems; it does not permit skipping the
  actual curve geometry. Formal fibres and canonical-model
  identification refer to BL85, Thuillier, and other foundations. Its
  Remark 4.2(2) cites Abbes for formal algebraization.
\item
  GH B.7 and B.8 were read as complete printed proof units. They use
  reduction/unit recognition and global formal-chart constructions.
  Proposition 12 independently replaces only the vertical node/unit and
  residual-degree calculation under its stated actual cover hypotheses.
\item
  KRZB Section 2.3 and Theorem 2.6 were read completely in their
  arbitrary complete algebraically closed nontrivial-field domain. The
  graph divisor is negative outgoing slope. Its footnote explicitly
  identifies the analytic extension beyond BPR's algebraic-function
  statement. No silent analytic extension is imported here. Its theorem
  excludes a singleton skeleton.
\item
  BE's actual projective model-to-pure-FS theorem 5.14 and nef-model
  recognition 7.8/7.9 were read in full. The latter requires an actual
  projective model and actual ample generic bundle. Lemma 6.3 supplies a
  finite norming \textbf{set} of Shilov points, not a single atom or a
  uniform individual lower bound. Reduced-fibre cofinality 4.19--4.20
  generates its own formal/reduction dependencies. GM 4.11/4.12 supplies
  ample algebraic domination only after the carrier is in its actual
  algebraic domain.
\item
  GM Appendix A.5 and A.7 were read fully. A.5 recognizes maximal
  residual points using affinoid dimension and formal embedding
  foundations. A.7 is limited to algebraic-completion extensions; A.8
  gives the warning for arbitrary extensions. Theorem 10 applies its
  elementary invariant calculation afresh on the actual generic curve,
  before the algebraic-completion step.
\item
  GJKM's {[}FK88{]} is \textbf{Eberhard Freitag and Reinhardt Kiehl},
  \emph{Étale cohomology and the Weil conjecture} (1988), as its
  bibliography confirms. The source's Proposition I.10.3.2 citation
  under that label is preserved as a source citation, not identified
  with another book. Fujiwara--Kato arXiv:1308.4734v5 is a separately
  read alternative primary. All new references use these distinct
  identities. The historical native191 misidentification remains
  historical input data only.
\item
  Fujiwara--Kato's whole Chapter 0 §9.2(a) proof and 8.5.21--25 were
  read. A rank-one complete valuation ring is \(a\)-adically complete
  and \(a\)-torsion-free: the valuation topology and \(a\)-adic topology
  have the same cofinal neighborhoods, and it is a domain. Corollary
  9.2.7 and Theorem 8.5.25(2) therefore have the intended base domain;
  their finiteness proof uses Raynaud--Gruson or its stated alternative
  division argument. These are not Noetherian/discrete hypotheses. Their
  complete generated foundations have not all been audited here.
\item
  The entire printed projective GFGA argument I.10.2(c), including
  10.2.4--6 and the final presentation argument, was read. It
  establishes existence on completed projective space using coherence,
  uniform graded-layer vanishing, comparison/full-faithfulness, and
  algebraization of finite presentations. The general carving proof and
  its Lemma 10.2.2 are abbreviated. I.10.3.2 explicitly leaves the
  adapted EGA III 5.4.5 algebraization argument to the reader. Reading
  its statement and reduction is not reading that missing proof.
  II.9.4--9.5's printed proofs treat the stronger t.u. adhesive base
  case, and generate comparison, formal coherent lift, and GFGA
  dependencies. Proposition 13 keeps the remaining
  section-identification premise explicit.
\end{itemize}

The supplied suggesting route through closed formal projective
embeddings is reasonable data to investigate. The preceding lifting
proof verifies the repeated-section part once whole mod-\(a\) generation
and vanishing are known. It does not establish those premises, finite
presentation/coherence of every required ideal, algebraization of the
carrier, or all compatibilities of the actual generic bundle. We do not
replace the omitted proof with an appeal to its statement.

\section{Remaining exact curve and higher-dimensional
obligations}\label{remaining-exact-curve-and-higher-dimensional-obligations}

R1--R5 from native191 are not all discharged. In the curve sector, the
remaining chain is: actual semistable domination of every required old
bundle/model in the non-Noetherian domain; a global node cover and
compatible refinements; whole mod-\(a\) ampleness/vanishing or another
verified actual same-bundle recognition; proper generic carrier and
bundle identification; genuine raw component-degree curvature
coefficients; complementary-ball frames and all terminal-point recovery;
and an actual finite Galois orbit for the downstairs norming reference.
Proposition 12 handles the explicit local Cartier and degree
calculation. Proposition 13 handles section lifting and FS recognition
with its stated premises. Neither turns an abstract graph distribution
into an actual metric without them.

The generic-incidence theorem constructs actual points on a common
auxiliary complete base under its initial residual-coordinate
hypotheses, and preserves every original value. It does not yet verify
those initial hypotheses for every finite reference Shilov set. The next
missing step is an actual common semistable curve core and domination
for the already actual restricted model/FS metrics, with the raw
slope/curvature formula. A finite-core voltage estimate would then bound
the reference Shilov values uniformly. Retain all graph loops:
refinements of a fixed full core can add trees and subdivide its edges,
and voltages extend constantly along attached trees. An abstract graph
estimate without that actual realization does not supply \(M_0\). No
general intrinsic integrability-to-finite-dual theorem is claimed here
without that bound.

Even a complete unconditional curve equation would leave the original
dimension-\(\geq2\) existence and every-solution all-point uniqueness
requirements. O and C already supply the operator and comparison in the
granted domain. T supplies the trivial boundary, and Q the split
projective line. For the remaining original equation there is still no
proved actual normalized-potential compactness or a replacement, no
complete envelope/orthogonality and both-sign outer differentiation
chain where that route needs it, no actual attainment for every
prescribed intrinsic target, and no every-solution recovery at every
point in arbitrary dimension. Corollary 6's attainment concerns an
already known output, and Corollary 5's coercivity concerns a still
conditional reference estimate. They do not establish the required
antecedents. The exact original problem remains open.

\section{Reproduction and evidence
limits}\label{reproduction-and-evidence-limits}

\texttt{verify\_packet.py} reads only the exact authorized manifest
files, verifies their bindings, and extracts the six full grants and
clarifications from \texttt{input.json}. \texttt{finite\_checks.py} uses
standard-library exact fractions and \(\mathbb Q(\sqrt2)\); it checks
the displayed interpolation prototype and 81 specified positive-bend
variants, finite Bernstein/scaling grids, raw curve normalization, and
nilpotence examples. Its conclusions are exhaustive only for those
finite grids. It checks no formal model, field-general psh theorem,
actual generic curve, or equation existence. Raw curve energy constants
are \(V^2/2\), as included in the exact checks.

\texttt{build\_deliverable.py} verifies inputs, reruns the checks,
generates the standalone TeX with existing pandoc, and compiles the PDF
twice with existing pdfLaTeX and shell escape disabled. It preserves
command arguments, return codes, output logs, environment versions, and
artifact hashes. Rerun in this directory:

\begin{Shaded}
\begin{Highlighting}[]
\ExtensionTok{python3}\NormalTok{ verify\_packet.py}
\ExtensionTok{python3}\NormalTok{ finite\_checks.py}
\ExtensionTok{python3}\NormalTok{ build\_deliverable.py}
\end{Highlighting}
\end{Shaded}

The permitted BL85 scan renders are in \texttt{scan-checks/}; their
generation and actual inspected printed pages are recorded in the
ledger. The primary PDFs remain at the exact read-only packet paths.
Public Stacks response snapshots are retained locally. No installation
or supervised resource-heavy computation was needed; no such compute
service was exposed, and none was launched locally. No protected editor,
other campaign stream, account/environment setting, private active
output, raw database, contact with another person, or GitHub mutation
was accessed or changed. Successful compilation and hash checks are
observations, not mathematical verification. All new mathematical
results remain unverified drafts.

\section{References and precise primary
identities}\label{references-and-precise-primary-identities}

O217/D214/N220/T171/Q222/C218: the six complete supplied accepted caller
bodies and all clarifications, preserved locally from
\texttt{input.json}.

Bosch--Lütkebohmert, \emph{Stable reduction and uniformization of
abelian varieties I}, Math. Ann. 270 (1985), 349--379, DOI
10.1007/BF01473432. Exact scanned primary:
\url{https://gdz.sub.uni-goettingen.de/download/pdf/PPN235181684_0270/LOG_0034.pdf}.

Gubler--Hertel, \emph{Local heights of toric varieties over
non-Archimedean fields}, literal arXiv:1512.06574v2:
\url{https://arxiv.org/pdf/1512.06574v2}.

Baker--Payne--Rabinoff, \emph{On the structure of nonarchimedean
analytic curves}, literal arXiv:1404.0279v1:
\url{https://arxiv.org/pdf/1404.0279v1}.

Katz--Rabinoff--Zureick-Brown, \emph{Uniform bounds for the number of
rational points on curves of small Mordell--Weil rank}, literal
arXiv:1504.00694v3: \url{https://arxiv.org/pdf/1504.00694v3}.

Boucksom--Eriksson, \emph{Spaces of norms, determinant of cohomology and
Fekete points in non-Archimedean geometry}, literal arXiv:1805.01016v3:
\url{https://arxiv.org/pdf/1805.01016v3}.

Gubler--Martin, \emph{On Zhang's semipositive metrics}, Documenta
Mathematica 24 (2019), 331--372:
\url{https://ems.press/content/serial-article-files/26466}.

Gubler--Jell--Künnemann--Martin, \emph{Continuity of plurisubharmonic
envelopes in non-Archimedean geometry and test ideals}, literal
arXiv:1712.00980v3: \url{https://arxiv.org/pdf/1712.00980v3}.

Freitag--Kiehl, \emph{Étale cohomology and the Weil conjecture},
Springer, 1988: the {[}FK88{]} bibliography identity in GJKM; its
original book proof was not read here.

Fujiwara--Kato, \emph{Foundations of Rigid Geometry I}, literal
arXiv:1308.4734v5, 28 February 2017, distinct from {[}FK88{]}:
\url{https://arxiv.org/pdf/1308.4734v5}.

Stacks Project, complete printed units in generic smooth Bertini
\url{https://stacks.math.columbia.edu/tag/0FD4}, generic irreducibility
\url{https://stacks.math.columbia.edu/tag/0G4C}, and the field lemmas
identified in the source ledger. These are primary data, not caller
grants. The incidence and Koszul arguments above state their actual
domains and additional premises explicitly.

\end{document}
